\subsection{Plane trigonometry.}

In this section, we choose a generator w of the algebra A(Φ) such that \(w^{2} \in A\) . Such a generator is determined up to a homothety (no. 1, Remark 2), hence the element \(w^{2}\) of A, which we denote by δ, is determined modulo the multiplicative subgroup (A*)² of squares of nonzero elements of A.

Remark. — When -1 belongs to the class mod. (A*)² in question, one generally chooses w such that \(w^{2} = -1\) , which determines it up to sign. When this class contains 1 but not -1, one generally chooses w such that \(w^{2} = 1\) , which again determines it up to sign.

This being so, every element \(\nu\) of \(S^{+}\) can be written, in one and only one way, in the form

\[
\nu = c _ {w} (\nu) + s _ {w} (\nu) \varpi\tag{1}
\]

where \(c_w(v)\) , \(s_w(v)\) belong to A; the element \(s_w(v)/c_w(v)\) of the projective field \(\tilde{\mathbf{A}}\) (chap.~II, 2 \(^{e}\) ed., App. III, no. 5) is denoted \(t_w(v)\) ; it depends only on the class of \(v\) mod. H; thus \(t_w\) defines, by passage to the quotient, a map from S \(^{+}\) /H into the projective field \(\tilde{\mathbf{A}}\) , which we shall still denote by \(t_w\) by abuse of language. We shall often write \(c\) , \(s\) and \(t\) instead of \(c_w\) , \(s_w\) and \(t_w\) . We have \(c_{-w} = c_w\) , \(s_{-w} = -s_w\) and \(t_{-w} = -t_w\) .

PROPOSITION 4. — a) When \(w^2 = \delta\) is not a square in A (that is, when E contains no isotropic line), the map t from S⁺/H into \(\tilde{\mathbf{A}}\) is a bijection.

\begin{enumerate}
\def\labelenumi{\alph{enumi})}
\setcounter{enumi}{1}
\item
  When δ is the square of an element γ of A, the map t is a bijection from S⁺/H onto \(\tilde{\mathbf{A}}\) minus the elements 1/γ and -1/γ.
\item
  Writing \(\mathrm{S}^{+}/\mathrm{H}\) additively, we have , for \(\varphi, \varphi'\) in \(\mathrm{S}^{+}/\mathrm{H}\)
\end{enumerate}

\[
t (\varphi + \varphi^ {\prime}) = (t (\varphi) + t (\varphi^ {\prime})) / (1 + \delta t (\varphi) t (\varphi^ {\prime}))\tag{2}
\]

when \(t(\varphi)\) and \(t(\varphi')\) are finite and \(1 + t(\varphi)t(\varphi')\) is \(\neq 0\) .

Indeed, as \(S^{+}/H\) is a set of lines (with 0 removed) of \(A(\Phi)\) considered as a vector plane over A, t is injective. On the other hand, for an element \(a + bw\) ( \(a \in A\) , \(b \in A\) ) of \(A(\Phi)\) to be a direct similarity, it is necessary and sufficient that it be invertible, that is, that one have \(N(a + bw) = a^{2} - \delta b^{2} \neq 0\) , or again \((b/a)^{2} \neq 1/\delta\) ; this proves the surjectivity assertions in a) and b). Finally, the product of the similarities \(1 + t(\varphi)w\) and \(1 + t(\varphi')w\) is the similarity \(1 + \delta t(\varphi)t(\varphi') + (t(\varphi) + t(\varphi'))w\) , which proves c).

PROPOSITION 5. --- Let us write O \(^{+}\) additively. For every pair of elements θ, θ' of O \(^{+}\) we have

\begin{enumerate}
\def\labelenumi{(\arabic{enumi})}
\setcounter{enumi}{2}
\tightlist
\item
\end{enumerate}

\[
c (\theta) ^ {2} - \delta s (\theta) ^ {2} = 1\tag{4}
\]

\[
c (\theta + \theta^ {\prime}) = c (\theta) c (\theta^ {\prime}) + \delta s (\theta) s (\theta^ {\prime})\tag{5}
\]

\[
s (\theta + \theta^ {\prime}) = s (\theta) c (\theta^ {\prime}) + c (\theta) s (\theta^ {\prime}).
\]

The relation (3) indeed expresses that \(\mathrm{N}(c(\theta)+s(\theta)\omega)=1\) To prove (4) and (5), it suffices to compute, in \(\mathrm{A}(\Phi)\) the product of the rotations \(c(\theta)+s(\theta)\omega\) and \(c(\theta')+s(\theta')\omega\) .

PROPOSITION 6. — Let d be the isomorphism from S+/H onto O+ defined in Prop. 3. For every element φ of S+/H such that t = t(φ) is finite, one has

\[
s (d (\varphi)) = 2 t / (1 - \delta t ^ {2}), \quad c (d (\varphi)) = (1 + \delta t ^ {2}) / (1 - \delta t ^ {2}).\tag{6}
\]

Indeed \(\varphi\) is the class mod. H of the similarity \(1 + tw\) , and the rotation \(d(\varphi)\) is therefore \((1 + tw)^2/\mathrm{N}(1 + tw) = (1 + \delta t^2 + 2tw)/(1 - \delta t^2)\) (prop. 3), which proves (6).

COROLLARY. --- For every element θ of O⁺ such that \(t(\theta)\) is finite, we have

\[
s (2 \theta) = 2 t (\theta) / (1 - \delta t (\theta) ^ {2}), \quad c (2 \theta) = (1 + \delta t (\theta) ^ {2}) / (1 - \delta t (\theta) ^ {2}). \tag {7}
\]

For every element \(\varphi\) of \(S^{+}/H\) such that \(t(\varphi)\) is finite and \(1 + \delta t(\varphi)^{2} \neq 0\) , one has

\[
t (2 \varphi) = 2 t (\varphi) / (1 + \delta t (\varphi) ^ {2}).\tag{8}
\]

Indeed this follows immediately from prop. 6 and the corollary of prop. 3.

Remark. --- Formulas (6) remain true for \(t = \infty\) provided that the rational functions appearing on the right-hand side are replaced by their canonical extensions to the projective field. \(\tilde{\mathbf{A}}\) (chap.~II, 2 \(^{\text{e}}\) ed., App. III, no. 5); indeed, if \(t = \infty\) , \(\varphi\) is the class of \(w\) , and we have \(d(\varphi) = -1\) , \(s(d(\varphi)) = 0\) and \(c(d(\varphi)) = -1\) ; these are indeed the values taken by the canonical extensions of the right-hand sides for \(t = \infty\) The same holds for (7) when \(t(\theta) = \infty\) , and for (8) when \(t(\varphi) = \infty\) or that \(1 + \delta t(\varphi)^{2} = 0\) . Similarly, formula (2) remains true when only one of the elements \(t(\varphi)\) , \(t(\varphi')\) , \(t(\varphi)\) for example, is infinite, provided that its right-hand side is regarded as a rational function of \(t(\varphi)\) alone: indeed, the product of the similarities \(1 + t(\varphi')w\) and \(w\) is \(\delta t(\varphi') + w\) , while the value taken by the canonical extension of the right-hand side of (2) is \(1/\delta t(\varphi')\) . Finally, when \(t(\varphi)\) and \(t(\varphi')\) are finite and one has \(1 + \delta t(\varphi)t(\varphi') = 0\) , one has \(t(\varphi) + t(\varphi') = 0\) (otherwise \(t(\varphi)^{2}\) would be equal to \(1/\delta\) , which is impossible (prop. 4)); one may therefore agree that the value of the right-hand side of (2) is \(\infty\) , and this value is indeed that of the left-hand side. When \(t(\varphi)\) and \(t(\varphi')\) are both infinite, the right-hand side of (2) is not defined.

